\chapter{Naive CA Theorem}

	\noindent
	I will put it directly here, without proof. It is so simple, that only saying it, you almost prove it. In our conversation it took you SECONDS to agree the concept.
	\\\\
	
	\noindent
	Being r a complete relation:\\
	$r: A \longrightarrow B $\\
	1) If we can FIND one PACK, not EMPTY, per each element of the Domain of r...\\
	2) If all those chosen PACKS, are DISJOINT all, between them...\\\\
	$\Longrightarrow Card( A ) \ngtr Card ( B ) $
	\\\\
	
	\noindent
	When ALL PACKS have cardinal 1, we have exactly the case when r is an injective function. Naive CA Th. is a generalization, for the injectivity practical use, for relations that are not functions.
	\\\\
	
	\noindent
	The game of Naive CA Th, consist in beginning with an original r relation, and DELETING R-pairs, from it, if necessary, until obtain a $r^{\prime}$ relation, that meets the conditions of Naive CA Th. We will proove our point always using a smaller relation, or equal, in quantity of R-pairs. THEN we will say the original r relation follows the NAIVE CA Th, so we can say the conclussion.
	\\\\
	
	\noindent
	This idea comes from my partner. Defining it using two relations. THIS IS WHY I create separate R-pairs per each image of the element of the Domain. We could play with the relations, deleting R-pairs, but what we left, after all the proccess... WAS ALWAYS, PREVIOUSLY, THERE. But:\\
	$ (a, \{2,4,6\}) \neq (a, \{4,6\}) $\\
	The R-pair of the right, could not be in the original r relation, and it could create the sensation that we are creating a NEW RELATION, adding stuff to adapt it to new propositions. When we are not adding R-pairs. We are going TO QUIT R-pairs. :D. Anyway... $r^{\prime}$ relation has clear conditions to be an equivalent of the practical use of injectivity. Focus on it if you are a little lost.
	\\\\
	
	\noindent
	The quantity of images we can use in each PACK has no limit. They don't need to be the same quantity in each PACK. They only need to exist, be not empty, and be DISJOINT. We will change the old word ``DIFFERENT'' for the word ``DISJOINT''. And our PACKs can have cardinal 1, 3, 4, millions... or directly infinite cardinal. THAT is going to be our case: PACKs with infinite cardinal... and they can not have a singular element in common. Between them all.
	\\\\
	
	\noindent
	I said to you in our conversation that this is a more general technic. That is why I am talking about useless stuff, sorry (for you). You gave me the meaning. So I will not need to explain DI (Inverse Diagonalization) where we play a lot quitting R-pairs from ``flja\_abstract''. A LOT !!. In our case, we will not quit R-pairs from relations $r_{\theta_{k}}$ (oriented more to TPI). You were convinced taking the meaning directly from a little resume of TPI technic. I have a document of 60 pages JUST for the definition of TPI (and properties). But you got it quickly... It happened to me before: someone hears the idea, seems legite... so they ask me for details, thinking the mistake is in the construction. Once they saw the construction is even more solid than the meaning... they try to change their original idea, trying to deny TPI... DI is needed, because THEY ALL, when they tried to deny TPI... guaranteed DI is perfectly builded. And when you try to deny DI, you can not avoid to guarantee TPI. This game is not neccesary for you. I say it... just in case... I saw it before... more than once. Please, don't be angry with me... :D.
	\\\\
	
	\noindent
	This paragraph is for other possible readers. IF YOU... until this point, got angry thinking I have not proven Naive CA Th., STOP READING. This is not for you. I am not one of your students, this is not an exam. This is a multidisciplinar solution, WHERE YOU ONLY NEED TO PUT A LITTLE EFFORT. Using all your knowledge. If you can not think in a quick proof of Naive CA Th... this means you can not adapt out of your comfort zone. It has around three different formal proofs, one not so formal, and the original one, that is mine, in my unortodox way of doing things. But I can not decide which one is better: I AM NOT MATHEMATICIAN. If you love rigor, over truth, stop reading. This is not for you. You will be able to read it in the future, exactly in the language and format you like, certified by your high priests. If you are more open minded:\\\\
	
	\noindent
	\textless ** \textit{See unnecesary comment 0005, more about Naive CA Th} \textgreater
	\\\\
	
	
	
	
	
	



\newpage